\loai{Lí thuyết}
\begin{ex} %[1P5B3-1][Loai1]
	Khi làm nóng một lượng khí có thể tích không đổi thì
	\choice
	{số phân tử trong một đơn vị thể tích giảm tỉ lệ nghịch với nhiệt độ}
	{áp suất khí không đổi}
	{số phân tử trong một đơn vị thể tích tăng tỉ lệ thuận với nhiệt độ}
	{\True số phân tử trong một đơn vị thể tích không đổi}
	\loigiai{
		Khi làm nóng một lượng khí có thể tích không đổi thì số phân tử trong một đơn vị thể tích không đổi.
		\begin{note}
			Khi đun nóng đẳng tích thì ta chỉ làm các phân tử chuyển động hỗn loạn hơn chứ mật độ phân tử vẫn không đổi.
		\end{note}
	}
\end{ex}
\begin{ex} %[1P5B3-1][Loai1]
	Đặc điểm nào sau đây \textbf{không} phải là quá trình biến đổi đẳng tích của một khối khí lí tưởng?
	\choice
	{Khi áp suất giảm chứng tỏ khối khí lạnh đi}
	{\True Áp suất khối khí tỉ lệ nghịch với nhiệt độ}
	{Áp suất của khối khí phụ thuộc vào nhiệt độ}
	{Khi nhiệt độ tăng thì áp suất của khối khí tăng}
	\loigiai{
		$\dfrac{p}{T} = \text{hằng số} \Rightarrow$ Áp suất khối khí tỉ lệ thuận với nhiệt độ.
	}
\end{ex}
\loai{Hệ thức quá trình đẳng tích}
\begin{ex} %[1P5B3-1][Loai2]
	Biểu thức nào sau đây \textbf{không} đúng cho định luật Sác-lơ?
	\choice
	{$\dfrac{p_1}{p_2} = \dfrac{T_1}{T_2}$}
	{$\dfrac{p_1}{T_1} = \dfrac{p_2}{T_2}$}
	{\True $\dfrac{p_1}{p_2} = \dfrac{T_2}{T_1}$}
	{${p_1}{T_2} = {p_2}{T_1}$}
	\loigiai{
		Theo định luật Sác-lơ: $\dfrac{p_1}{T_1} = \dfrac{p_2}{T_2}\Rightarrow \dfrac{p_1}{p_2} = \dfrac{T_1}{T_2}$.
	}
	%<MyLT>
\end{ex}
\begin{ex}%[1P5Y3-1][Loai2]
	Hệ thức nào sau đây phù hợp với định luật Sác-lơ?
	\choice
	{$p\sim t$}
	{$\dfrac{p}{t} = \text{hằng số}$}
	{\True $\dfrac{p_1}{T_1} = \dfrac{p_3}{T_3}$}
	{$p_1t_1=p_2t_2$}
	\loigiai{
		Ta chỉ có $p\sim T$ chứ \textbf{không} có $p\sim t$. 	
	}
\end{ex}
\begin{ex} %[1P5Y3-1][Loai2]
	Hệ thức nào sau đây phù hợp với định luật Sác-lơ?	
	\choice
	{$p \sim \dfrac{1}{T}$}
	{${p_1}{T_1} = {p_2}{T_2}$}
	{$\dfrac{p_1}{T_2} = \dfrac{p_2}{T_1}$}
	{\True $\dfrac{p}{T} =$ hằng số}
	
	\loigiai{
		Hệ thức phù hợp với định luật Sác-lơ: $\dfrac{p}{T} = \text{hằng số}$.
	}
\end{ex}
\begin{ex} %[1P5Y3-1][Loai2]
	Khi làm nóng một lượng khí có thể tích không đổi, gọi $n_0$ là số phân tử trong một đơn vị thể tích, p là áp suất, T là nhiệt độ tuyệt đối. Tỉ số nào sau đây không đổi?
	
	\choice
	{$\dfrac{n_0}{p}$}
	{\True $\dfrac{p}{T}$}
	{$\dfrac{n_0}{pT}$}
	{$\dfrac{n_0}{T}$}
	
	\loigiai{
		Quá trình đẳng tích nên $\dfrac{p}{T}= \text{hằng số}$.
	}
	%<MyLT>
\end{ex}
\begin{ex}%[1P5Y3-1][Loai2]
	Trong các hệ thức sau đây, hệ thức nào \textbf{không} phù hợp với định luật Sác-lơ?
	\choice
	{$p\sim T$}
	{\True $p\sim t$}
	{$\dfrac{p}{T} = \text{hằng số}$}
	{$\dfrac{p_1}{T_1} = \dfrac{p_2}{T_2}$}
	\loigiai{
		Ta chỉ có $p\sim T$ chứ \textbf{không} có $p\sim t$. 	}
\end{ex}
\loai{Đồ thị p-T}
\begin{ex} %[1P5Y3-1][Loai3]
	Trong hệ tọa độ (p, T), đường biểu diễn nào sau đây là đường đẳng tích?
	
	\choice
	{\True Đường thẳng kéo dài qua gốc tọa độ}
	{Đường thẳng không đi qua gốc tọa độ}
	{Đường hypebol}
	{Đường thẳng cắt trục p tại điểm $p = p_0$}
	
	\loigiai{
		$\dfrac{p}{T} = \text{hằng số}$ $\Rightarrow$ Đường thẳng kéo dài qua gốc tọa độ.
	}
\end{ex}
\begin{ex} %[1P5Y3-1][Loai3]
	Trong đồ thị (p, T), đường đẳng tích là
	\choice
	{đường thẳng vuông góc với trục p}
	{\True đường thẳng đi qua O}
	{đường thẳng vuông góc với trục T}
	{đường hyperbol}
	
	\loigiai{
		Trên đồ thị (p, T), đường đẳng tích là đường thẳng đi qua O.
	}
\end{ex}
\begin{ex}%[1P5K3-1][Loai3]
	immini[thm]{Cho đồ thị của áp suất theo nhiệt độ của hai khối khí A và B
		có thể tích không đổi như hình vẽ. Nhận xét nào sau đây là \textbf{sai}?\choice
		{Hai đường biểu diễn đều cắt trục hoành tại điểm $-273\;\doC$}
		{Khi $t = 0\;\doC$, áp suất của khối khí A lớn hơn áp suất của khối khí B}
		{Áp suất của khối khí A luôn lớn hơn áp suất của khối khí B tại mọi nhiệt độ}
		{\True Khi tăng nhiệt độ, áp suất của khối khí B tăng nhanh hơn áp suất của khối khí A}}{\Vehinh{1}{
			\draw[->] (-2.5,0)--(3,0)node[below]{$t$}; \draw[->] (0,0)--(0,3)node[right]{p};
			\draw[dashed] (-2,0)--++(30:1);
			\draw (-2,0)++(30:1)--++(30:4)node[right]{$A$};
			\draw[dashed] (-2,0)--++(10:1);
			\draw (-2,0)++(10:1)--++(10:4)node[right]{$B$};
	}}
	
	
	\loigiai{
		Độ dốc của $p_A$ lớn hơn độ dốc của $p_B$ nên áp suất A tăng nhanh hơn B.
	}
	%<MyLT>
\end{ex}
\begin{ex} %[1P5B3-1][Loai3]
	immini[thm]{Một khối khí lí tưởng thực hiện quá trình đẳng tích ở hai thể tích khác nhau được biểu diễn trên hình vẽ. Quan hệ giữa $V_1$ và $V_2$ là
		\choice
		{${V_1}={V_2}$}
		{${V_1}<{V_2}$}
		{\True ${V_1}>{V_2}$}
		{không so sánh được}
	}{\begin{tikzpicture}[scale=1,thick,>=stealth']
			\draw[->] (0,0)--(3,0)node[below]{$T$}; \draw[->] (0,0)--(0,3)node[left]{p};
			\draw[dashed] (0,0)--++(60:1);
			\draw (0,0)++(60:1)--++(60:2)node[right]{$V_2$};
			\draw[dashed] (0,0)--++(30:1);
			\draw (0,0)++(30:1)--++(30:2)node[right]{$V_1$};
	\end{tikzpicture}}
	
	\loigiai{
		immini[thm]{\begin{itemize}
				\item Từ đồ thị $\Rightarrow$ ứng với nhiệt độ $T_1$: $P_1<P_2$. 
				\item Mà theo định luật Bôi-lơ - Ma-ri-ốt ta có $P_1V_1 = P_2V_2 \Rightarrow V_1 > V_2$.
		\end{itemize}}{\begin{tikzpicture}[scale=1,thick,>=stealth']
				\coordinate (o) at (0,0);
				\node at (o) [below left]{\normalsize{O}};
				\draw [->](o)--++(0:4.5)coordinate(t)node[below]{\normalsize{T}};
				\draw [->](o)--++(90:4)coordinate(p)node[left]{\normalsize{p}};
				\draw [dashed](o)--++(30:1)coordinate(v2);
				\draw (v2) -- ($(v2)!3cm!180:(o)$)node[right]{\normalsize{\bth{V_{1}}}};
				\draw [dashed](o)--++(60:1)coordinate(v1);
				\draw (v1) -- ($(v1)!3cm!180:(o)$)node[above]{\normalsize{\bth{V_{2}}}};
				\path (o) --++(0:2) node[below left]{$T_1$};
				\path(o)--++(0:1.5)coordinate(1);
				\draw[dashed](1)--++(90:0.85)coordinate(a)--++(90:1.7)coordinate(b)--++(180:1.5)coordinate(p1)node[left]{\normalsize{\bth{p_{2}}}};
				\draw[dashed](a)--++(180:1.5)coordinate(p2)node[left]{\normalsize{\bth{p_{1}}}};
				
			\end{tikzpicture}
		}
	}
\end{ex}
\begin{ex}%[1P5B3-1][Loai3]
	immini[thm]{Cho đồ thị $p - T$ biểu diễn hai đường đẳng tích của cùng một khối khí xác định như hình vẽ. Đáp án nào sau đây biểu diễn đúng mối quan hệ về thể tích của hai đường?
		\choice
		{$V_1 > V_2$}
		{\True $V_1 < V_2$}
		{$V_1 = V_2$}
		{$V_1 \ge V_2$}
	}{\Vehinh{1}{
			\draw[->] (0,0)--(3,0)node[below]{$T$}; \draw[->] (0,0)--(0,3)node[left]{p};
			\draw[dashed] (0,0)--++(60:1);
			\draw (0,0)++(60:1)--++(60:2)node[right]{$V_1$};
			\draw[dashed] (0,0)--++(20:1);
			\draw (0,0)++(20:1)--++(20:2)node[right]{$V_2$};
	}}
	
	\loigiai{
		\begin{itemize}
			\item Kẻ đường song song OT cắt các đường đẳng tích $V_1$ và $V_2$ tại các điểm M và N.
			\item Từ M và N kẻ các đường song song trục Op, cắt trục OT tại các điểm có hoành độ $T_1$ và $T_2$.
			\item Vì p không đổi, $T_1<T_2$ nên $V_1<V_2$.
		\end{itemize}
	}
\end{ex}

\loai{Đồ thị T- p}
\begin{ex} %[1P5K3-1][Loai4]
	immini[thm]{Cùng một khối lượng khí của một chất khí đựng trong 3 bình kín có thể tích khác nhau, đồ thị thay đổi áp suất theo nhiệt độ của ba khối khí ở ba bình được mô tả như hình vẽ. Quan hệ về thể tích của ba bình đó là
		\choice
		{${V_3} = {V_2} = {V_1}$}
		{\True $V_3<V_2<V_1$ }
		{${V_3} > {V_2} > {V_1}$}
		{${V_3} \ge {V_2} \ge {V_1}$}
	}{\begin{tikzpicture}[scale=1,thick,>=stealth']
			\draw[->] (0,0)--(3,0)node[below]{$p$}; \draw[->] (0,0)--(0,3)node[left]{T};
			\draw[dashed] (0,0)--++(45:1);
			\draw (0,0)++(45:1)--++(45:2)node[right]{$V_2$};
			\draw[dashed] (0,0)--++(60:1);
			\draw (0,0)++(60:1)--++(60:2)node[right]{$V_1$};
			\draw[dashed] (0,0)--++(30:1);
			\draw (0,0)++(30:1)--++(30:2)node[right]{$V_3$};
	\end{tikzpicture}} 
	
	\loigiai{
		immini[thm]{\begin{itemize}
				\item Ứng với nhiệt độ $T_1$, $p_1 < p_2 < p_3$. 
				\item Mà theo định luật Bôi-lơ - Ma-ri-ốt: ${p_1}{V_1} = {p_2}{V_2} = {p_3}{V_3} \Rightarrow {V_1} > {V_2} > {V_3}$.
		\end{itemize}}{\begin{tikzpicture}[scale=1,thick,>=stealth']
				\draw[->] (0,0)--(3,0)node[below]{$p$}; \draw[->] (0,0)--(0,3)node[left]{T};
				\draw[dashed] (0,0)--++(45:1);
				\draw (0,0)++(45:1)--++(45:2)node[right]{$V_2$};
				\draw[dashed] (0,0)--++(60:1);
				\draw (0,0)++(60:1)--++(60:2)node[right]{$V_1$};
				\draw[dashed] (0,0)--++(30:1);
				\draw (0,0)++(30:1)--++(30:2)node[right]{$V_3$};
				\draw[dashed] (0,1.25)node[left]{$T_1$}--(2.15,1.25)--(2.15,0)node[below]{$p_3$}(1.25,0)node[below]{$p_2$}--(1.25,1.25)(0.73,0)node[below]{$p_1$}--(0.73,1.25);
		\end{tikzpicture}}
	}
	%<MyLT>
\end{ex}

\loai{Đồ thị V-T}
\begin{ex} %[1P5B3-1][Loai5]
	Trên đồ thị (V, T), đường đẳng tích là
	\choice
	{đường thẳng vuông góc với trục T}
	{đường thẳng có phương đi qua gốc tọa độ}
	{\True đường thẳng vuông góc với trục V}
	{đường hyperbol}
	\loigiai{
		Trên đồ thị (V, T), đường đẳng tích là đường thẳng vuông góc với trục V.
	}
\end{ex}
\loai{Đổi hệ trục}
\begin{ex} %[1P5B3-1][Loai6]
	Đồ thị nào sau đây thể hiện định luật Sác-lơ?	
	\choice
	{\begin{tikzpicture}[scale=0.8,thick,>=stealth']
			\draw[->] (-1,0)--(3,0)node[below]{$T$}; \draw[->] (0,0)--(0,3)node[right]{$p$};
			\draw[dashed] (-0.5,0)--++(30:1);
			\draw (-0.5,0)++(30:1)--++(30:2.5);
			\node[below](0,0){$O$};
	\end{tikzpicture}}
	{\True \begin{tikzpicture}[scale=0.8,thick,>=stealth']
			\draw[->] (0,0)--(3,0)node[below]{$V$}; \draw[->] (0,0)--(0,3)node[right]{$p$};
			\draw (1,0)--++(90:2.5);
			\node[below](0,0){$O$};
	\end{tikzpicture}}
	{\begin{tikzpicture}[scale=0.8,thick,>=stealth']
			\draw[->] (0,0)--(3,0)node[below]{$t$}; \draw[->] (0,0)--(0,3)node[right]{$p$};
			\draw (0,0)--++(45:3);
			\node[below](0,0){$O$};
	\end{tikzpicture}}
	{\begin{tikzpicture}[scale=0.8,thick,>=stealth']
			\draw[->] (0,0)--(3,0)node[below]{$V$}; \draw[->] (0,0)--(0,3)node[right]{$p$};
			\draw (0,0)--++(45:3);
			\node[below](0,0){$O$};
	\end{tikzpicture}}
	
	\loigiai{
		Theo định luật Sác-lơ: V = hằng số, $\dfrac{P}{T}=$ hằng số.
	}
\end{ex}
\begin{ex} %[1P5B3-1][Loai6]
	immini[thm]{Một khối khí thay đổi trạng thái như đồ thị. Đồ thị trên được vẽ sang hệ trục p – V là hình nào trong các hình sau đây?
		\choice
		{\begin{tikzpicture}[scale=0.7,thick,>=stealth']
				\tkzDefPoints{0/0/o, 4/0/t, 0/4/v, 3/3/1, 1.5/3/2, 1.5/1.5/3}
				\tkzDrawSegments[->](o,t o,v)
				\tkzDrawSegments[arrow data={0.5}{stealth}](3,2 2,1)
				\draw[dashed](1)--(3,0)node[below]{$2V_0$};
				\draw[dashed](3)--(1.5,0)node[below]{$V_0$};
				\draw[dashed](3)--(0,1.5)node[left]{$p_0$};
				\draw[dashed](2)--(0,3)node[left]{$2p_0$};
				\tkzLabelPoint[above](t){$V$}
				\tkzLabelPoint[above](1){$(3)$}
				\tkzLabelPoint[above](2){$(2)$}
				\tkzLabelPoint[below left](3){$(1)$}
				\tkzLabelPoint[right](v){$p$}
				\tkzLabelPoint[below left](o){$O$}
		\end{tikzpicture}}
		{\True \begin{tikzpicture}[scale=0.7,thick,>=stealth']
				\tkzDefPoints{0/0/o, 4/0/v, 0/4/p, 1.5/1.5/1, 1.5/3/2, 3/1.5/3}
				\tkzDrawSegments[->](o,p o,v)
				\tkzDrawSegments[arrow data={0.5}{stealth}](1,2)
				\draw[arrow data={0.5}{stealth}](2) to [bend right =30](3);
				\draw[dashed](1)--(1.5,0)node[below]{$V_0$};
				\draw[dashed](3)--(3,0)node[below]{$2V_0$};
				\draw[dashed](3)--(0,1.5)node[left]{$p_0$};
				\draw[dashed](2)--(0,3)node[left]{$2p_0$};
				\tkzLabelPoint[above](v){$V$}
				\tkzLabelPoint[above](3){$(3)$}
				\tkzLabelPoint[above](2){$(2)$}
				\tkzLabelPoint[below left](1){$(1)$}
				\tkzLabelPoint[right](p){$p$}
				\tkzLabelPoint[below left](o){$O$}
		\end{tikzpicture}}
		{\begin{tikzpicture}[scale=0.7,thick,>=stealth']
				\tkzDefPoints{0/0/o, 4/0/v, 0/4/p, 1.5/1.5/1, 1.5/3/2, 3/1.5/3}
				\tkzDrawSegments[->](o,p o,v)
				\tkzDrawSegments[arrow data={0.5}{stealth}](1,2)
				\tkzDrawSegments[arrow data={0.5}{stealth}](2,3)
				\draw[dashed](1)--(1.5,0)node[below]{$V_0$};
				\draw[dashed](3)--(3,0)node[below]{$2V_0$};
				\draw[dashed](3)--(0,1.5)node[left]{$p_0$};
				\draw[dashed](2)--(0,3)node[left]{$2p_0$};
				\tkzLabelPoint[above](v){$V$}
				\tkzLabelPoint[above](3){$(3)$}
				\tkzLabelPoint[above](2){$(2)$}
				\tkzLabelPoint[below left](1){$(1)$}
				\tkzLabelPoint[right](p){$p$}
				\tkzLabelPoint[below left](o){$O$}
		\end{tikzpicture}}
		{\begin{tikzpicture}[scale=0.7,thick,>=stealth']
				\tkzDefPoints{0/0/o, 4/0/v, 0/4/p, 1.5/1.5/1, 3/3/2, 1.5/3/3}
				\tkzDrawSegments[->](o,p o,v)
				\tkzDrawSegments[arrow data={0.5}{stealth}](1,2 2,3)
				\draw[dashed](3)--(1.5,0)node[below]{$V_0$};
				\draw[dashed](2)--(3,0)node[below]{$2V_0$};
				\draw[dashed](1)--(0,1.5)node[left]{$p_0$};
				\draw[dashed](2)--(0,3)node[left]{$2p_0$};
				\tkzLabelPoint[above](v){$V$}
				\tkzLabelPoint[above](3){$(3)$}
				\tkzLabelPoint[above](2){$(2)$}
				\tkzLabelPoint[below left](1){$(1)$}
				\tkzLabelPoint[right](p){$p$}
				\tkzLabelPoint[below left](o){$O$}
			\end{tikzpicture}
		}
	}{\begin{tikzpicture}[scale=0.8,thick,>=stealth']
			\tkzDefPoints{0/0/o, 4/3/b, 4/1.5/a, 2/1.5/c}
			\draw[->] (0,0)--(5,0)node[below]{$T$}; 
			\draw[->] (0,0)--(0,3.5)node[left]{$p$};
			\draw[dashed](0,1.5)--(c)--(2,0)(0,3)--(b)(b)--(4,0)(o)--(c)--(a);
			\tkzDrawSegments[blue,very thick,arrow data={0.5}{stealth}](b,a c,b)
			\node at (a)[right]{\small(3)};
			\node at (b)[right]{\small(2)};
			\node at (c)[above left]{\small (1)};
			\node at (0,3)[left]{\small $2p_{0}$};
			\node at (0,1.5)[left]{\small $p_{0}$};
			\node at (2,0)[below]{\small $T_{0}$};
			\node at (0.5,0.8)[]{\small $V_{0}$};
		\end{tikzpicture}
	} 
	
	\loigiai{
		Từ (1) đến (2) là quá trình đẳng tích, từ (2) đến (3), nhiệt độ không đổi, áp suất giảm $\Rightarrow$ là quá trình dãn đẳng nhiệt.
	}
\end{ex}
\begin{ex} %[1P5B3-1][Loai6]
	immini[thm]{Một khối khí lý tưởng thực hiện quá trình được biểu diễn như trên hình vẽ. Đồ thị nào cũng biểu diễn đúng quá trình trên?
		\choice
		{\begin{tikzpicture}[scale=0.8,thick,>=stealth']
				\draw[->] (0,0)--(3.5,0)node[below]{$T$}; 
				\draw[->] (0,0)--(0,3.5)node[right]{$V$};
				\draw[dashed](0,2)node[left]{$V_1=V_2$}--(1,2)node[above,inner sep=1pt]{\small(1)}--(3,2)node[above,fill=white,inner sep=1pt]{\small (2)};
				\draw[dashed](1,2)--(1,0)node[below=3pt,inner sep=1pt]{\small$T_1$}(3,2)--(3,0)node[below=3pt,fill=white,inner sep=1pt]{\small$T_2$};
				\draw[undirected,blue, ultra thick](3,2)--(1,2);
		\end{tikzpicture}}
		{\begin{tikzpicture}[scale=0.8,thick,>=stealth']
				\draw[->] (0,0)--(4,0)node[below]{$T$}; 
				\draw[->] (0,0)--(0,3.5)node[right]{V};
				\draw[dashed](0,1)node[left]{$V_1$}--(1,1)node[above left,inner sep=1pt]{\small(1)}--(1,0)node[below]{$T_1$};
				\draw[dashed](0,3)node[left]{$V_2$}--(3,3)node[right,inner sep=1pt]{\small(2)}--(3,0)node[below]{$T_2$};
				\draw[directed,blue, ultra thick](3,3)--(1,1);
		\end{tikzpicture}}
		{\True \begin{tikzpicture}[scale=0.8,thick,>=stealth']
				\draw[->] (0,0)--(3.5,0)node[below]{$p$}; 
				\draw[->] (0,0)--(0,3.5)node[right]{$V$};
				\draw[dashed](0,2)node[left]{$V_1=V_2$}--(1,2)node[above,inner sep=1pt]{\small(2)}--(3,2)node[above,inner sep=1pt]{\small (1)};
				\draw[dashed](1,2)--(1,0)node[below=3pt,inner sep=1pt]{\small$p_2$}(3,2)--(3,0)node[below=3pt,fill=white,inner sep=1pt]{\small$p_1$};
				\draw[undirected,blue, ultra thick](1,2)--(3,2);
		\end{tikzpicture}}
		{\begin{tikzpicture}[scale=0.8,thick,>=stealth']
				\draw[->] (0,0)--(4,0)node[below]{$T$}; 
				\draw[->] (0,0)--(0,3.5)node[right]{p};
				\draw[dashed](0,1)node[left]{$p_1$}--(1,1)node[above left,inner sep=1pt]{\small(1)}--(1,0)node[below]{$T_1$};
				\draw[dashed](0,3)node[left]{$p_2$}--(3,3)node[right,inner sep=1pt]{\small(2)}--(3,0)node[below]{$T_2$};
				\draw[directed,blue, ultra thick](1,1)--(3,3);
		\end{tikzpicture}}
	}{\begin{tikzpicture}[scale=0.8,thick,>=stealth']
			\draw[->] (0,0)--(3.5,0)node[below]{$T$}; 
			\draw[->] (0,0)--(0,3.5)node[right]{$V$};
			\draw[dashed](0,2)node[left]{$V_1=V_2$}--(1,2)node[above,inner sep=1pt]{\small(2)}--(3,2)node[above,fill=white,inner sep=1pt]{\small (1)};
			\draw[dashed](1,2)--(1,0)node[below=3pt,inner sep=1pt]{\small$T_2$}(3,2)--(3,0)node[below=3pt,fill=white,inner sep=1pt]{\small$T_1$};
			\draw[undirected,blue, ultra thick](1,2)--(3,2);
	\end{tikzpicture}} 
	
	
	\loigiai{
		Từ đồ thị $\Rightarrow$ quá trình từ (1) đến (2) là quá trình đẳng tích.
	}
\end{ex}

